% Source: https://artofproblemsolving.com/community/c3h1513891_amc8_p22
% Stable identity: aops-topic-1513891-q01
Place the rectangle with $E=(0,0)$, $F=(3,0)$, $A=(3,4)$, $D=(0,4)$; then $C=(1,4)$ and $B=(2,4)$.

\textbf{The point $P$ is determined by the figure.} It is the intersection of the lines $EB$ and $FC$, so it is not an extra assumption:
\[
EB:\ y=2x,\qquad FC:\ y=-2(x-3).
\]
Setting $2x=-2x+6$ gives $x=\dfrac{3}{2}$ and $y=3$, so $P=\left(\dfrac{3}{2},3\right)$. In particular $P$ is $1$ unit below side $DA$; this is computed, not guessed.

\begin{asy}
usepackage("amsmath");
size(4.5cm);
defaultpen(fontsize(10pt));
pair E=(0,0),F=(3,0),A=(3,4),D=(0,4),C=(1,4),B=(2,4);
pair P=extension(E,B,F,C);
fill(E--C--P--cycle,rgb(0.75,0.85,1));fill(F--B--P--cycle,rgb(0.75,0.85,1));
draw(E--F--A--D--cycle,linewidth(1));
draw(E--B--F,linewidth(0.8));draw(F--C--E,linewidth(0.8));
dot(P,red);
label("$A(3,4)$",A,NE);label("$B$",B,N);label("$C$",C,N);label("$D(0,4)$",D,NW);label("$E(0,0)$",E,SW);label("$F(3,0)$",F,SE);
label("$P$",P,3.5*S,red);
label("\mbox{}",(min(currentpicture,user=true)+max(currentpicture,user=true))/2,black); // keeps the text after the figure black when LaTeX typesets the labels inline
\end{asy}

\textbf{Area of one wing.} By the shoelace formula, triangle $ECP$ with $E=(0,0)$, $C=(1,4)$, $P=\left(\dfrac{3}{2},3\right)$ has area
\[
\dfrac{1}{2}\left|1\cdot3-4\cdot\dfrac{3}{2}\right|=\dfrac{1}{2}\cdot3=\dfrac{3}{2} .
\]
By symmetry about the line $x=\dfrac{3}{2}$, triangle $FBP$ has the same area. The bat wings have area $2\cdot\dfrac{3}{2}=3$, answer $\boxed{\textbf{(C) }3}$.
